\scanpage{159}
% archive foliation: 158
Catégorie $\mathcal C$, sous-catégorie $\mathcal P_{\unsure{1}}$ (pleine), morphismes ``non ramifiés''

\uline{$1^\circ$} Dans $\mathcal C$, les sommes finies existent, ainsi que les produits fibrés.

\uline{$4^\circ$} Dans $\mathcal C$, le \deleted{produit}\inserted{composé} de deux morphismes non ramifiés \inserted{les morphismes identiques}, ou l'image inverse d'un morphisme non ramifié, sont non ramifiés.

\uline{$5^\circ$} \deleted{\unsure{les}} $f:X\to Y$ est non ramifié $\Longleftrightarrow$ les $f|X_i$ ($X_i$ les comp. connexes) \deleted{\unsure{les sommes $\coprod(f_i:X_i\to Y)$}} non ramifiés.

\uline{$2^\circ$} Dans $\mathcal C$, il y a des composantes connexes.

\uline{$3^\circ$} Les \deleted{pts} $\xi\in\mathcal P$ sont connexes
(donc $\operatorname{Hom}(\xi,X\amalg Y)=\operatorname{Hom}(\xi,X)\amalg\operatorname{Hom}(\xi,Y)$]

\uline{$6^\circ$} Pour tout $f:X\to Y$ non ramifié, \inserted{$Y$ connexe}, et $\xi\in\mathcal P$,
$\operatorname{Hom}(\xi,X)\to\operatorname{Hom}(\xi,Y)$ est \uline{surjectif}.

\uline{$7^\circ$} Si \deleted{$X$} et $Y$ sont connexes, et la fibre de
$\operatorname{Hom}(\xi,X)\to\operatorname{Hom}(\xi,Y)$ est réduite à un élément, alors $f$ est un \uline{isom}
(on suppose $f$ non ramifié, $\xi\in\mathcal P$)

\uline{$8^\circ$} Si $f:X\to Y$ non ramifié, $Y$ connexe, alors \deleted{pour} toute \inserted{toute} section $s:Y\to X$ est un \uline{isom.} de $Y$ sur une composante connexe de $X$.

\uline{$9^\circ$} Soit \deleted{$G$ \inserted{groupe fini} opère dans $X$ ``sans \unsure{points}''}
$f:X\to Y$ \uline{non ramifié}, et $X$ connexe, \unsure{$\xi\in$} $e\in\operatorname{Hom}(\xi,Y)=Y^\xi$, tel que $G=\operatorname{Aut}_Y(X)$ opère transitivement dans $(f^\xi)^{-1}(e)$.
Alors pour tout sous-groupe $H$ de $G$, $X/H$ existe \inserted{dans $\mathcal C$}, et est un revêtement non ramifié de $Y$.
\deleted{et $X\to\unsure{X/H}$ \unsure{est non ramifié}}
et deux élts de $\operatorname{Hom}_{\unsure{Y}}(\xi,X)$ au-dessus de $e$ ayant $=$ image par $X\to X/H$ sont congrus mod $H$.

Sous ces conditions, pour $x\in\operatorname{Hom}(\xi,X)$ on définit $\pi_1(X,x)$.
Il ne change pas si on remplace $\xi$ par $\xi'$ tel que $\xi'\to\xi$ soit donné.
C'est un foncteur \uline{covariant} en $(X,x)$

\diagram{\begin{tikzcd}
\xi \arrow[r,"x"] \arrow[d] & X \arrow[d] \\
\xi' \arrow[r,"x'"] & X'
\end{tikzcd}}
définissant $\pi_1(X,x)\mathbin{\unsure{-|\to}}\pi_1(X',x')$.

\uline{Etc.}
